% GENERATED FILE. Edit canonical lessons, then run npm run generate:books.
% canonical-source-hash: fnv1a64:42f5a3978dd7b542
% canonical-lessons: JA-C242-kare, JA-C242-kanojo, JA-C242-koibito, JA-C242-shinyuu, JA-C242-tsukiau

\chapter{Partners and Best Friends}
\label{ch:ja-partners-and-best-friends}
\hlchaptermodality{242}

\begin{chapteropening}
\textbf{By the end of this chapter:} \emph{I can ask and answer questions about friends and partners.}

Complete the last lesson of chapter 242: I can ask and answer questions about friends and partners.
\end{chapteropening}

\section[\texorpdfstring{\ja{かれ} (kare)}{かれ (kare)}]{\ja{かれ} (kare) — he; a boyfriend}

\label{lesson:JA-C242-kare}



\begin{quote}
Before the new one: say the Japanese for to sign up, then the Japanese for to do.
\end{quote}

\subsection*{You'll want to know: \ja{かれ}}
\textbf{\ja{かれ}} — \emph{kare} — \textquotedblleft{}he; a boyfriend\textquotedblright{}.

\begin{quote}
\textbf{\ja{かれは} \ja{がくせいですか。}} — \emph{kare wa gakusei desu ka} — \textquotedblleft{}Is he a student?\textquotedblright{}
\end{quote}

\begin{cousinweb}
\textbf{\ja{かれ}} (\emph{kare}) is a native Japanese word (wago), inherited rather than borrowed from Chinese; its kanji was attached later as a kun'yomi (JMdict, KANJIDIC2).
\end{cousinweb}

\begin{grammarlens}[title={What you've built}]
One more word for everyday life at A2.
\end{grammarlens}

\subsection*{Your turn}
\begin{itemize}
\raggedright
  \item \emph{Say it:} \emph{kare}
  \item \emph{Say it:} \emph{kare}, once more
  \item \emph{Say it:} \emph{kare wa gakusei desu ka} — Is he a student?
  \item \emph{Recall:} say the Japanese for to sign up, then the Japanese for to do, then say \emph{kare} again
\end{itemize}

\begin{tcolorbox}[breakable,colback=teal!4,colframe=teal!35!black,title={Before you move on}]
What does \ja{かれ} mean? (\textquotedblleft{}He; a boyfriend\textquotedblright{}.) Say it once more.
\end{tcolorbox}
\section[\texorpdfstring{\ja{かのじょ} (kanojo)}{かのじょ (kanojo)}]{\ja{かのじょ} (kanojo) — she; a girlfriend}

\label{lesson:JA-C242-kanojo}



\begin{quote}
Before the new one: say the Japanese for to do, then the Japanese for he.
\end{quote}

\subsection*{You'll want to know: \ja{かのじょ}}
\textbf{\ja{かのじょ}} — \emph{kanojo} — \textquotedblleft{}she; a girlfriend\textquotedblright{}.

\begin{quote}
\textbf{\ja{かのじょは} \ja{いま} \ja{きょうしつに} \ja{いません。}} — \emph{kanojo wa ima kyōshitsu ni imasen} — \textquotedblleft{}She is not in the classroom now.\textquotedblright{}
\end{quote}

\begin{cousinweb}
\textbf{\ja{かのじょ}} (\emph{kanojo}) is a mixed word: native \emph{kano}, \textquotedblleft{}that\textquotedblright{}, joined to Sino-Japanese \emph{jo}, \textquotedblleft{}woman\textquotedblright{} (Middle Chinese \emph{nrjoX}), borrowed from Chinese (JMdict, KANJIDIC2; Middle Chinese in Baxter's transcription).
\end{cousinweb}

\begin{grammarlens}[title={What you've built}]
One more word for everyday life at A2.
\end{grammarlens}

\subsection*{Your turn}
\begin{itemize}
\raggedright
  \item \emph{Say it:} \emph{kanojo}
  \item \emph{Say it:} \emph{kanojo}, once more
  \item \emph{Say it:} \emph{kanojo wa ima kyōshitsu ni imasen} — She is not in the classroom now.
  \item \emph{Recall:} say the Japanese for to do, then the Japanese for he, then say \emph{kanojo} again
\end{itemize}

\begin{tcolorbox}[breakable,colback=teal!4,colframe=teal!35!black,title={Before you move on}]
What does \ja{かのじょ} mean? (\textquotedblleft{}She; a girlfriend\textquotedblright{}.) Say it once more.
\end{tcolorbox}
\section[\texorpdfstring{\ja{こいびと} (koibito)}{こいびと (koibito)}]{\ja{こいびと} (koibito) — a sweetheart; a lover}

\label{lesson:JA-C242-koibito}



\begin{quote}
Before the new one: say the Japanese for he, then the Japanese for she.
\end{quote}

\subsection*{You'll want to know: \ja{こいびと}}
\textbf{\ja{こいびと}} — \emph{koibito} — \textquotedblleft{}a sweetheart; a lover\textquotedblright{}.

\begin{quote}
\textbf{\ja{あにには} \ja{こいびとが} \ja{いません。}} — \emph{ani ni wa koibito ga imasen} — \textquotedblleft{}My elder brother does not have a sweetheart.\textquotedblright{}
\end{quote}

\begin{cousinweb}
\textbf{\ja{こいびと}} (\emph{koibito}) is a native Japanese compound (wago), inherited rather than borrowed: it is built from \emph{koi}, \textquotedblleft{}love\textquotedblright{}, and \emph{hito}, \textquotedblleft{}person\textquotedblright{}, with \emph{h} voiced to \emph{b} (JMdict, KANJIDIC2).
\end{cousinweb}

\begin{grammarlens}[title={What you've built}]
One more word for everyday life at A2.
\end{grammarlens}

\subsection*{Your turn}
\begin{itemize}
\raggedright
  \item \emph{Say it:} \emph{koibito}
  \item \emph{Say it:} \emph{koibito}, once more
  \item \emph{Say it:} \emph{ani ni wa koibito ga imasen} — My elder brother does not have a sweetheart.
  \item \emph{Recall:} say the Japanese for he, then the Japanese for she, then say \emph{koibito} again
\end{itemize}

\begin{tcolorbox}[breakable,colback=teal!4,colframe=teal!35!black,title={Before you move on}]
What does \ja{こいびと} mean? (\textquotedblleft{}A sweetheart; a lover\textquotedblright{}.) Say it once more.
\end{tcolorbox}
\section[\texorpdfstring{\ja{しんゆう} (shin'yū)}{しんゆう (shin'yū)}]{\ja{しんゆう} (shin'yū) — a close friend; a best friend}

\label{lesson:JA-C242-shinyuu}



\begin{quote}
Before the new one: say the Japanese for she, then the Japanese for a sweetheart.
\end{quote}

\subsection*{You'll want to know: \ja{しんゆう}}
\textbf{\ja{しんゆう}} — \emph{shin'yū} — \textquotedblleft{}a close friend; a best friend\textquotedblright{}.

\begin{quote}
\textbf{\ja{かれは} \ja{あにの} \ja{しんゆうですか。}} — \emph{kare wa ani no shin'yū desu ka} — \textquotedblleft{}Is he my elder brother's best friend?\textquotedblright{}
\end{quote}

\begin{cousinweb}
\textbf{\ja{しんゆう}} (\emph{shin'yū}) is Sino-Japanese (kango): two Chinese characters read with their on'yomi, readings borrowed from Chinese. They are \emph{shin}, \textquotedblleft{}intimate\textquotedblright{} (Middle Chinese \emph{tshin}) and \emph{yū}, \textquotedblleft{}friend\textquotedblright{} (\emph{hjuwX}) (KANJIDIC2; Middle Chinese in Baxter's transcription).
\end{cousinweb}

\begin{grammarlens}[title={What you've built}]
One more word for everyday life at A2.
\end{grammarlens}

\subsection*{Your turn}
\begin{itemize}
\raggedright
  \item \emph{Say it:} \emph{shin'yū}
  \item \emph{Say it:} \emph{shin'yū}, once more
  \item \emph{Say it:} \emph{kare wa ani no shin'yū desu ka} — Is he my elder brother's best friend?
  \item \emph{Recall:} say the Japanese for she, then the Japanese for a sweetheart, then say \emph{shin'yū} again
\end{itemize}

\begin{tcolorbox}[breakable,colback=teal!4,colframe=teal!35!black,title={Before you move on}]
What does \ja{しんゆう} mean? (\textquotedblleft{}A close friend; a best friend\textquotedblright{}.) Say it once more.
\end{tcolorbox}
\section[\texorpdfstring{\ja{つきあう} (tsukiau)}{つきあう (tsukiau)}]{\ja{つきあう} (tsukiau) — to go out with; to keep company with}

\label{lesson:JA-C242-tsukiau}



\begin{quote}
Before the new one: say the Japanese for a sweetheart, then the Japanese for a close friend.
\end{quote}

\subsection*{You'll want to know: \ja{つきあう}}
\textbf{\ja{つきあう}} — \emph{tsukiau} — \textquotedblleft{}to go out with; to keep company with\textquotedblright{}.

\begin{quote}
\textbf{\ja{あねは} \ja{かれと} \ja{つきあいますか。}} — \emph{ane wa kare to tsukiaimasu ka} — \textquotedblleft{}Will my elder sister go out with him?\textquotedblright{}
\end{quote}

Polite \textbf{\ja{つきあいます}}, \emph{tsukiaimasu}; a godan verb.

\begin{cousinweb}
\textbf{\ja{つきあう}} (\emph{tsukiau}) is a native Japanese compound (wago), inherited rather than borrowed: it is built from \emph{tsuki}, from \emph{tsuku}, \textquotedblleft{}to stick to\textquotedblright{}, and \emph{au}, \textquotedblleft{}to do together\textquotedblright{} (JMdict, KANJIDIC2).
\end{cousinweb}

\begin{grammarlens}[title={What you've built}]
One more word for everyday life at A2.
\end{grammarlens}

\subsection*{Your turn}
\begin{itemize}
\raggedright
  \item \emph{Say it:} \emph{tsukiau}
  \item \emph{Say it:} \emph{tsukiau}, once more
  \item \emph{Say it:} \emph{ane wa kare to tsukiaimasu ka} — Will my elder sister go out with him?
  \item \emph{Recall:} say the Japanese for a sweetheart, then the Japanese for a close friend, then say \emph{tsukiau} again
\end{itemize}

\begin{tcolorbox}[breakable,colback=teal!4,colframe=teal!35!black,title={Before you move on}]
What does \ja{つきあう} mean? (\textquotedblleft{}To go out with; to keep company with\textquotedblright{}.) Say it once more.
\end{tcolorbox}
